\chapter*{怎样阅读本册}
\addcontentsline{toc}{chapter}{怎样阅读本册}

阅读第三册时，不要先把注意力放在模型名和框架名上。模型名称会变，框架接口会变，硬件指令集也会扩展；但本地大模型推理的基本账本长期稳定：权重字节从哪里来，token 怎样变成 id，张量按什么步长读取，在哪个类型中累加，KV cache 写到哪里，CPU/GPU 后端各自承担哪些 kernel，怎样判断误差可接受，怎样证明时间数字对应正确工作。

如果你完全不懂编译原理，也可以从本册开始。先把“编译器”理解成一个严格的工程流水线：读懂输入，建立中间表示，检查不变量，在语义不变的前提下改写成更适合机器执行的形式。AI 编译器只是把这件事用在模型文件、张量、量化、推理图、后端 kernel 和运行时调度上。遇到 IR、verifier、lowering、pass、runtime 这些词，不要先背定义，而要问：它在本地推理引擎里检查了什么、生成了什么、节省了哪部分机器成本。

本册会频繁回扣第一册和第二册。第一册教你把 C++ 代码追到 ABI、汇编、对象布局和可信 benchmark；第二册教你把同一段机器工作放进现代 CPU、cache/TLB、SIMD、线程、NUMA、运行时和 I/O 边界里解释。第三册讲 AI 编译器和推理后端时，会把这两类能力当成工具，而不是重新孤立发明一套 AI 术语。

建议按四条线并行阅读。

\begin{enumerate}
  \item 架构线：模型资产、编译前端、IR 优化、后端、运行时分别负责什么，哪些事实不能跨层乱猜。
  \item 语义线：每个算子定义什么输入、输出和错误边界。
  \item 表示线：token、张量、权重、scale、zero point、activation buffer 和 KV cache 怎样存在内存里。
  \item 机器线：核心循环生成什么地址序列、使用哪些指令、触发哪些 cache、带宽、NUMA、PCIe 或显存问题。
  \item 后端线：同一个算子合同怎样分别落到 CPU SIMD/多线程 kernel 和 GPU kernel/stream 上。
  \item 证据线：reference、误差 oracle、反汇编、profiler、benchmark 和框架对标报告怎样组成闭环。
\end{enumerate}

若某个优化看起来很诱人，先问五件事：它是否保持算子语义，是否改变了机器账本中真正等待的部分，是否适用于 prefill、decode 或两者，C++20 工程实现是否有清楚模块边界和资源所有权，是否有测量证据证明它相对当前基线和现有框架缩小了差距。只要这些问题没有说清，就不要把它当成可靠优化。

\section*{固定主案例：MiniLLM-CPU}

第三册所有重要概念都应回到同一个主案例：\term{MiniLLM-CPU}{一个从 tiny fixture 到 7B shape 账本的本地推理引擎切片}。它有两个尺度。第一个尺度小到可以手算，用来训练正确性；第二个尺度接近真实 7B，用来训练容量、带宽和后端判断。

\begin{center}
\begin{tabular}{p{0.20\linewidth}p{0.35\linewidth}p{0.31\linewidth}}
\toprule
尺度 & 固定参数 & 用途 \\
\midrule
tiny fixture & \code{vocab=16, hidden=8, layers=2, heads=2, kv\_heads=1, head\_dim=4, seq=4} & 手算 shape、stride、RoPE、GQA、KV cache、logits golden \\
kernel fixture & \code{T=2, K=8, N=6} 与 \code{H=4096, O=4096, K=4096} & 对照 naive、packed、SIMD、tail、cache 和带宽上限 \\
7B shape & \code{L=32, H=4096, heads=32, kv\_heads=8, head\_dim=128, ctx=4096} & 估算权重字节、KV 容量、prefill/decode tokens/s 上限 \\
\bottomrule
\end{tabular}
\end{center}

tiny fixture 贯穿的 token 链路固定为：

\begin{lstlisting}[numbers=none,caption={MiniLLM-CPU 的固定 token trace}]
prompt ids
  -> embedding [T,H]
  -> RMSNorm [T,H]
  -> Q [T,heads,head_dim], K/V [T,kv_heads,head_dim]
  -> RoPE on Q/K
  -> KV cache append [layer,position,kv_head,head_dim]
  -> GQA attention output [T,H]
  -> MLP SwiGLU [T,H]
  -> lm_head logits [T,vocab]
  -> sampler next token
\end{lstlisting}

每一章都必须说明自己接在这条链路的哪里。讲张量，要给 shape、stride 和地址公式；讲量化，要给 byte layout、scale 和误差；讲 compiler，要给 IR dump、pass 前后差异和 plan；讲 backend，要给 baseline、反汇编、perf 或替代计数器；讲 serving，要给 queue、TTFT、TPOT、取消和拒绝。

\section*{每章最低密度}

第三册不接受只讲“应该有什么”的章节。每章至少要留下下面十类内容中的多数项，核心章节必须全部覆盖：

\begin{rubricbox}
\begin{enumerate}
  \item 一个具体问题，不用抽象口号开头。
  \item 一个 tiny 数值例子，能手算或能人工检查。
  \item 一个真实 shape 例子，至少说明 7B 级别的字节或延迟影响。
  \item 一段可编译或可直接迁移到 C++20 项目的核心代码。
  \item 一张地址流、内存布局、状态机或时间线图。
  \item 一张 benchmark、ledger 或 oracle 表。
  \item 一段 perf、反汇编、trace 或等价工具证据。
  \item 一个错误案例，例如 shape 错、tail 错、KV offset 错、scale 错、fallback 静默发生。
  \item 一个作业，要求读者改代码或复现实验。
  \item 一个和 MiniLLM-CPU artifact 的连接：fixture、命令、CSV、report、CI 或 release gate。
\end{enumerate}
\end{rubricbox}

\section*{学完第三册应能做什么}

学完第三册不等于可以胜任所有 AI infra 岗位，但应达到一个明确出口：

\begin{itemize}
  \item 能写出 tiny decoder 的 reference，并用 logits golden 定位 embedding、RMSNorm、QKV、RoPE、attention、KV cache、MLP 和 sampler 的错误。
  \item 能为 7B 级模型计算 fp16/int8/int4 权重字节、KV cache 容量、decode 最小读流和带宽 tokens/s 上限。
  \item 能实现一个 int8 GEMV/GEMM 路径：scalar baseline、packed layout、SIMD 或清晰后端接口、tail path、operator oracle 和 benchmark。
  \item 能设计 AI compiler 的最小 IR：TensorDesc、GraphNode、ShapeVerifier、MemoryPlan、BackendLowering、ExecutablePlan 和 pass verifier。
  \item 能把本地推理接成服务：admission、session、KV 预算、prefill/decode 调度、stream event、取消、backpressure、SLO trace。
  \item 能写一份工程证据报告：correctness、benchmark、perf/trace、工业框架对标、已知限制和下一步优化。
\end{itemize}

若这些做不到，说明只是读过 AI infra 名词；若这些能稳定完成，才说明第三册真正把第一册和第二册的系统能力用到了本地大模型推理上。
